% ===========================================================
%  第 5 章 线性方程组 —— 高等代数【深化】拓展框（子代理转录）
%  素材：北大《高等代数》第五版 第三章；樊启斌《高等代数典型问题与方法》第三章、§4.4
%  说明：与 deep_ch05_authored.tex 的 3 块（核空间 / 特解加平移 / 列空间有解判别）
%        不重复，只补 Cramer 边界、同解判定、基础解系构造、导出组、含参数讨论等面。
% ===========================================================

% ---------- Cramer 法则的适用边界（北大 三.6 印刷页 93–94；樊 3.2 印刷页 139）----------
\begin{fdeep}{Cramer 法则的两层用法：算唯一解，和"给一般方程组收尾"}
樊 3.2 的陈述：数域 $K$ 上含 $n$ 个未知量 $n$ 个方程的方程组，若系数行列式 $D\ne0$，则它有唯一解
\fml{x_j=\frac{D_j}{D}\qquad (j=1,2,\dots,n)}
其中 $D_j$ 是把 $D$ 的第 $j$ 列换成常数项 $b_1,\dots,b_n$ 得到的行列式。\textbf{边界很硬}：必须有 $n$ 个方程、$n$ 个未知量，且 $D\ne0$；$D=0$ 时它一个字也说不了；方程个数与未知量个数不等时，$D$ 根本不存在。

北大 三.6（印刷页 93–94）给出的是它的\textbf{理论用法}——用 Cramer 法则解一般方程组：设 $r(A)=r(\bar A)=r$，取 $A$ 的一个不为零的 $r$ 阶子式 $D$（不妨设位于左上角），则 $\bar A$ 的前 $r$ 行构成 $\bar A$ 的一个极大线性无关组，第 $r+1,\dots,s$ 行都可经它们线性表出，故\textbf{删去后 $s-r$ 个方程不改变解集}。于是：
\begin{itemize}[leftmargin=2.2em,itemsep=0.25em]
\item $r=n$ 时，直接对留下的 $r$ 个方程用 Cramer 法则，得唯一解；
\item $r<n$ 时，把含 $x_{r+1},\dots,x_n$ 的项移到等号右边当作已知量，剩下 $r$ 个方程的系数行列式恰为 $D\ne0$，于是\textbf{自由未知量每取一组值，都用 Cramer 法则得到唯一解}。
\end{itemize}
可见 Cramer 法则并不因 $D=0$ 而"失效"，它只是要求你先选好保留哪 $r$ 个方程、把哪些未知量挪到右边。
【反哺点】：服务考纲"会用克拉默法则"。考场计算只用第一层；理解结构用第二层。第二层顺手说清了基础解系为什么恰好 $n-r$ 个——正是被挪到右边当参数的那 $n-r$ 个未知量。
\end{fdeep}

% ---------- 齐次情形的 Cramer 推论与它的地盘（樊 3.2 印刷页 140；樊 3.3 印刷页 142）----------
\begin{fdeep}{齐次方程组：$D\ne0$ 与"只有零解"互为充要，但"$D$"只在方阵上有}
由 Cramer 法则立得：当 $b_1=\cdots=b_n=0$ 而 $D\ne0$ 时，$Ax=0$ 只有零解；取逆否，$Ax=0$ 有非零解则必有 $D=0$。方阵情形这两句其实是\textbf{充要}的，因为
\fml{D=0\iff r(A)<n\iff Ax=0\ \text{有非零解}}
但 $D$ 只在\textbf{方程个数等于未知量个数}时才存在。对一般的 $m\times n$ 系数矩阵，唯一判据是秩：
\fml{Ax=0\ \text{有非零解}\iff r(A)<n}
注意 $n$ 是\textbf{未知量个数}，与方程个数 $m$ 无关；特别地 $m<n$ 时必然有非零解（樊 3.3.2）。
实操中 $D=0$ 最常见的用法是\textbf{反解参数}：把题干条件写成"某齐次方程组有非零解"，凑出 $n$ 阶系数矩阵，由 $D=0$ 解出参数候选值。这只用到了必要性，\textbf{候选值必须逐个代回}：$D=0$ 只保证秩不满，题目往往还附加"基础解系恰含 $k$ 个向量""某给定向量是解"等条件，要再算秩来筛。
【反哺点】：服务考纲"齐次线性方程组有非零解的充分必要条件"。一句话记住："$D\ne0$"只在方阵上管用，一般情形一律换成 $r(A)<n$；而 $D=0$ 只负责筛参数，筛完必须代回验算。
\end{fdeep}

% ---------- 有解的七种等价说法（樊 3.4.1 印刷页 157）----------
\begin{fdeep}{$Ax=\beta$ 有解的七种等价说法：把"有解"翻译成子空间语言}
设 $A=(\alpha_1,\dots,\alpha_n)$ 为 $m\times n$ 矩阵，$\beta$ 为 $m$ 维列向量。下列说法\textbf{完全等价}（樊 3.4.1）：
\begin{enumerate}[leftmargin=2.2em,itemsep=0.2em]
\item 非齐次方程组 $Ax=\beta$（即 $\sum_{j=1}^{n}x_j\alpha_j=\beta$）有解；
\item $\beta$ 可由列向量组 $\alpha_1,\dots,\alpha_n$ 线性表示；
\item 向量组 $\alpha_1,\dots,\alpha_n$ 与 $\alpha_1,\dots,\alpha_n,\beta$ 等价；
\item 它们张成的子空间相同：$L(\alpha_1,\dots,\alpha_n)=L(\alpha_1,\dots,\alpha_n,\beta)$；
\item 两个张成子空间的维数相同，即 $\dim L(\alpha_1,\dots,\alpha_n)=\dim L(\alpha_1,\dots,\alpha_n,\beta)$；
\item 向量组的秩相等：$r(\alpha_1,\dots,\alpha_n)=r(\alpha_1,\dots,\alpha_n,\beta)$；
\item 矩阵的秩相等：$r(A)=r(A,\beta)$。
\end{enumerate}
同一套语言还能翻译"解的情形"：$r(A)=r(A,\beta)=n$ 时表示法\textbf{唯一}；$r(A)=r(A,\beta)<n$ 时能表示但表示法\textbf{不唯一}；$r(A)\ne r(A,\beta)$ 时\textbf{不能}表示。
【反哺点】：服务考纲"理解非齐次线性方程组有解的充分必要条件"。这七条里只有秩的两条是"能算的"，其余是"能懂的"——"有解"的几何身份就是\textbf{添上 $\beta$ 之后张成空间没有变大}。命题判断题最爱在"有解 / 可线性表示 / 张成空间相同 / 秩相等"之间做替换，认出它们是一回事，题就送分。
\end{fdeep}

% ---------- 无解的精确判别与平面几何对照（北大 三.5 印刷页 93、三.6 印刷页 98）----------
\begin{fdeep}{无解的精确形状：$r(\bar A)=r(A)+1$；以及它的平面几何对照}
增广矩阵 $\bar A$ 只比 $A$ 多一列，所以两个秩至多差 1：
\fml{r(A)\le r(\bar A)\le r(A)+1}
于是"无解"有一个\textbf{精确}写法（北大 三.5 印刷页 93）：
\fml{Ax=b\ \text{无解}\iff r(\bar A)=r(A)+1}
（相应地"有解"就是 $r(\bar A)=r(A)$。）讨论含参数方程组时这条特别省事：只要阶梯形里出现 $0\ 0\ \cdots\ 0\mid d$（$d\ne0$）这样一行，就直接断定无解，不必再回头核对两个秩的数值。

\textbf{几何对照}（北大 三.6 印刷页 98）以两个方程三个未知量为例，每个方程表示一张平面，$r(A)$ 只能是 1 或 2：
\begin{itemize}[leftmargin=2.2em,itemsep=0.25em]
\item $r(A)=1$、$r(\bar A)=1$：$A$ 的两行成比例、$\bar A$ 的两行也成比例，两平面\textbf{重合}，无穷多解；
\item $r(A)=1$、$r(\bar A)=2$：两平面\textbf{平行而不重合}，无解；
\item $r(A)=2$：此时必有 $r(\bar A)=2$，两平面\textbf{不平行因而必定相交}，无穷多解。
\end{itemize}
判别式在这里不是算出来的，而是"两张平面相对位置"的代数翻译。
【反哺点】：服务考纲"理解非齐次线性方程组有解的充分必要条件"。$r(\bar A)=r(A)+1$ 去掉了"$r(A)<r(\bar A)$ 到底小多少"的含糊；几何对照则让"解的个数"类选择题可以不计算就排除选项（三张平面的情形同理，$r(A)=r(\bar A)=2$ 对应交于一条直线）。
\end{fdeep}

% ---------- 初等变换保同解的证明思路（北大 三.1 印刷页 71–72）----------
\begin{fdeep}{初等变换为什么保同解：一个"互相包含"的证明模板}
三种变换称为线性方程组的\textbf{初等变换}：① 用一非零数乘某一方程；② 把一个方程的倍数加到另一个方程；③ 互换两个方程的位置。消元的过程就是反复施行初等变换的过程。核心结论：\textbf{初等变换总是把方程组变成同解方程组}（北大 三.1）。
证明思路（以第②种为例）：设把第 2 个方程的 $k$ 倍加到第 1 个方程得新方程组。若 $(c_1,\dots,c_n)$ 是原方程组的任一解，则它自动满足新方程组中未改动的后 $s-1$ 个方程；又因原前两式成立，
\fmll{a_{11}c_1+\cdots+a_{1n}c_n=b_1,\qquad a_{21}c_1+\cdots+a_{2n}c_n=b_2,}
把第二式两边乘 $k$ 再加到第一式，正得 $(a_{11}+ka_{21})c_1+\cdots+(a_{1n}+ka_{2n})c_n=b_1+kb_2$，即新方程组的第一个方程。故原解都是新解；对称地可证新解都是原解。两个解集互相包含，因而不相上下——这就是同解。两种变换（乘非零数、互换位置）的证明同法可做。
【反哺点】：服务考纲"掌握用初等行变换求解线性方程组的方法"。它回答了"凭什么可以对增广矩阵随便做行变换"——每一步都保持解集不变。凡遇"变换后解集是否改变"的辨析题，就用"两个解集互相包含"这两句话作答。
\end{fdeep}

% ---------- 同解方程组的判定（樊 3.3 印刷页 143–144）----------
\begin{fdeep}{同解方程组：秩相等只是必要条件，不是充分条件}
两个方程组\textbf{同解} $\iff$ 它们有相同的解集合。设 $A,B$ 都是 $m\times n$ 矩阵。若 $Ax=0$ 与 $Bx=0$ 同解，则二者的解空间相同，于是
\fml{n-r(A)=n-r(B)\ \Longrightarrow\ r(A)=r(B)}
\textbf{但反过来不成立}。樊 3.3 给出的反例：
\fml{\begin{pmatrix}1&2\\0&0\end{pmatrix}\begin{pmatrix}x_1\\x_2\end{pmatrix}=0,\qquad \begin{pmatrix}2&1\\0&0\end{pmatrix}\begin{pmatrix}x_1\\x_2\end{pmatrix}=0}
两个系数矩阵的秩都是 1，但前者要求 $x_1=-2x_2$（$x_2$ 任意），后者要求 $x_2=-2x_1$（$x_1$ 任意），显然不同解。可见\textbf{秩相等只说明两个解空间的维数相同，并不说明是同一个子空间}。
正确的判定是：$Ax=0$ 与 $Bx=0$ 同解 $\iff$ 两者的\textbf{行向量组等价}（可以互相线性表出）——因为每个方程都是一次线性约束，"同解"意味着两套约束能互相推出。特别地，若 $A$ 可经初等行变换化为 $B$，则两方程组同解。
【反哺点】：服务考纲"理解齐次线性方程组的基础解系、通解及解空间的概念"。命题判断题里"同解"与"秩相等"最易混淆：要判同解，要么比较两个解空间（解出通解比一比），要么检验行向量组能否互相表出；只比秩一定不够。
\end{fdeep}

% ---------- 基础解系标准求法的构造性证明（北大 三.6 印刷页 96；樊 3.3 印刷页 142–143）----------
\begin{fdeep}{"自由未知量取单位向量"为什么必然得到一组基}
定理（北大 三.6）：齐次方程组有非零解时必有基础解系，且基础解系所含解的个数等于 $n-r$，其中 $r$ 为系数矩阵的秩。\textbf{这个定理的证明就是求法本身}。
把行最简形写回方程组，以 $x_{r+1},\dots,x_n$ 为自由未知量。任给它们一组值就唯一确定一个解；换句话说，\textbf{自由未知量的值一样，两个解就完全一样}。现在取 $n-r$ 组单位向量当自由未知量的值：
\fml{(1,0,\dots,0),\ (0,1,\dots,0),\ \dots,\ (0,0,\dots,1)}
得到 $n-r$ 个解（樊 3.3 印刷页 143 给出了它们的显式形状）。
\begin{itemize}[leftmargin=2.2em,itemsep=0.25em]
\item \textbf{线性无关}：若 $k_1\eta_1+\cdots+k_{n-r}\eta_{n-r}=0$，比较\textbf{最后 $n-r$ 个分量}（它们恰是 $k_1,\dots,k_{n-r}$），立得 $k_1=\cdots=k_{n-r}=0$；
\item \textbf{能表出任一解}：设 $\eta$ 是任一解，则 $c_{r+1}\eta_1+\cdots+c_n\eta_{n-r}$（$c_{r+1},\dots,c_n$ 是 $\eta$ 的自由未知量取值）也是解，且与 $\eta$ 的最后 $n-r$ 个分量相同，故两个解完全相同。
\end{itemize}
【反哺点】：服务考纲"掌握基础解系和通解的求法"。"自由未知量依次取单位向量、比较最后 $n-r$ 个分量"就是考场的标准模板；这一招还能反过来用于\textbf{快速判断给定的一组解向量是否线性无关}（只看末 $n-r$ 个分量即可）。
\end{fdeep}

% ---------- 基础解系不唯一的等价刻画（北大 三.6 印刷页 97；樊 3.3 印刷页 142）----------
\begin{fdeep}{基础解系不唯一：任意 $n-r$ 个线性无关的解都是一组基}
樊 3.3.1 的两条定义性事实：齐次方程组 $Ax=0$ 的全部解构成 $K^n$ 的一个子空间，称为\textbf{解空间}；非零解空间的一个\textbf{基}称为基础解系；且 $\dim S=n-r(A)$。
由定义立刻看出：\textbf{任何一个与某个基础解系等价的线性无关向量组，仍是一个基础解系}（北大 三.6 印刷页 97）。理由有两条：其一，"能线性表出全部解"这个性质在等价替换下保持不变；其二，等价的两个线性无关组所含向量个数相同。
由此得到两条实用推论：
\begin{itemize}[leftmargin=2.2em,itemsep=0.25em]
\item \textbf{基础解系不唯一，但所含向量个数唯一}，恒为 $n-r$。这解释了为什么不同解法（换自由未知量、换消元顺序）写出的基础解系长相不同却都对——它们张成的是同一个子空间。
\item \textbf{$Ax=0$ 的任意 $n-r$ 个线性无关的解向量必构成基础解系}：个数已经"够多"，线性无关就说明它们撑满了整个解空间。
\end{itemize}
【反哺点】：服务考纲"理解齐次线性方程组的基础解系、通解及解空间的概念"。"任意 $n-r$ 个线性无关的解就是基础解系"是选择题里出现频率极高的一句话；它把"验证基础解系"从"验证能表出所有解"（难）降级成"数个数 + 查线性无关"（易）。
\end{fdeep}

% ---------- 导出组与唯一性判据（北大 三.6 印刷页 97–98；樊 3.4.3 印刷页 158）----------
\begin{fdeep}{导出组：非齐次解集的两条运算法则与"唯一解"判据}
把 $Ax=b$ 的常数项全部换成 $0$ 得到的齐次方程组 $Ax=0$，称为 $Ax=b$ 的\textbf{导出组}（樊 3.4.3、北大 三.6）。两者解之间有两条运算法则（北大 三.6 印刷页 97）：
\begin{itemize}[leftmargin=2.2em,itemsep=0.25em]
\item $Ax=b$ 的\textbf{两个解之差}是其导出组 $Ax=0$ 的解；
\item $Ax=b$ 的\textbf{一个解加上导出组的一个解}，仍是 $Ax=b$ 的解。
\end{itemize}
由第一条，取任一特解 $\gamma_0$，任一解 $\gamma$ 都可写成 $\gamma=\gamma_0+(\gamma-\gamma_0)$，其中 $\gamma-\gamma_0$ 是导出组的解，于是全部解为 $\gamma_0+\eta$（$\eta$ 取遍导出组的解）——这正是通解 $\gamma_0+k_1\xi_1+\cdots+k_{n-r}\xi_{n-r}$ 的来源。
\textbf{推论（常考）}：在 $Ax=b$ \textbf{有解}的前提下，
\fml{Ax=b\ \text{的解唯一}\iff \text{导出组}\ Ax=0\ \text{只有零解}}
证明用第一条即可：若 $Ax=b$ 有两个不同解，其差是导出组的非零解；反之若导出组有非零解，把它加到任一解上就得到另一个解。换成秩语言：方阵情形就是 $\det A\ne0$，一般情形就是 $r(A)=n$（列满秩）。
【反哺点】：服务考纲"理解非齐次线性方程组解的结构及通解的概念"。"非齐次解唯一 $\iff$ 导出组只有零解"比"$r(A)=r(\bar A)=n$"更好用：含参数题里只要判出齐次部分只有零解，就能立刻断定唯一，不必同时盯两个秩。
\end{fdeep}

% ---------- $AB=O$ 与解空间（樊 4.4.3 印刷页 221）----------
\begin{fdeep}{$AB=O$ 为什么给出 $r(A)+r(B)\le n$：看 $B$ 的列落在谁里面}
樊 4.4.3（印刷页 221）的秩不等式：设 $A$ 是 $m\times n$ 矩阵、$B$ 是 $n\times s$ 矩阵，若 $AB=O$，则
\fml{r(A)+r(B)\le n}
考研辅导里这条常被死记，其实\textbf{用方程组一句话就说清了}。把 $B$ 按列分块 $B=(\beta_1,\dots,\beta_s)$，则 $AB=O$ 等价于
\fml{A\beta_j=0,\qquad j=1,2,\dots,s}
即 \textbf{$B$ 的每一列都是齐次方程组 $Ax=0$ 的解}——$B$ 的列向量组整个落在解空间（核空间）里。于是
\fml{r(B)=\dim L(\beta_1,\dots,\beta_s)\le \dim N(A)=n-r(A)}
移项即得 $r(A)+r(B)\le n$。\textbf{等号何时成立}：当 $B$ 的列向量组恰好是 $Ax=0$ 的一组基时，$r(B)=n-r(A)$，等号成立。所以"$AB=O$ 且 $r(A)+r(B)=n$"这句话里其实藏着"$B$ 的列构成 $Ax=0$ 的基础解系"。
【反哺点】：服务考纲"理解齐次线性方程组的基础解系、通解及解空间的概念"。秩不等式里 $r(A)+r(B)\le n$（注意是 $A$ 的\textbf{列数} $n$，不是 $m$ 也不是 $s$）最容易记错，用"$B$ 的列都在解空间里"这一步就永远不会写错；反过来，题目一给 $AB=O$，就要立刻想到"$B$ 的列是 $Ax=0$ 的解"。
\end{fdeep}

% ---------- 含参数方程组的讨论路线（樊 3.3 印刷页 144、3.4 印刷页 157–158 例题归纳）----------
\begin{fdeep}{含参数方程组：先分流，再分情形}
由樊 §3.3、§3.4 的含参数例题（如例 3.32、例 3.55 等）可归纳出一条稳定的讨论路线：
\begin{enumerate}[leftmargin=2.2em,itemsep=0.25em]
\item \textbf{方程个数等于未知量个数 $n$}：先算系数行列式 $D$。$D\ne0$ 时由 Cramer 法则直接得\textbf{唯一解}（此时不必讨论）；再令 $D=0$ 解出参数候选值，逐个代回，对增广矩阵作初等行变换，比较 $r(A)$ 与 $r(\bar A)$。
\item \textbf{方程个数不等于未知量个数}（或参数藏在常数项里）：不算行列式，\textbf{直接对增广矩阵作初等行变换}，把含参数的表达式保留在矩阵中，找出"会使某个主元消失"的参数值，再按这些值分情形。
\item 两条省力的判别捷径：① 阶梯形里出现 $0\ 0\ \cdots\ 0\mid d$（$d\ne0$）$\Longrightarrow$ 无解（即 $r(\bar A)=r(A)+1$）；② 若已知方程组\textbf{有不止一个解}（例如题给三个互异的解向量），则其导出组有非零解；当系数矩阵是 $n$ 阶方阵时立刻得到 $\det A=0$。
\end{enumerate}
【反哺点】：服务考纲"非齐次线性方程组有解的充分必要条件"及含参数方程组的讨论。先把题目分流到第 1 类还是第 2 类，能省掉大量无用计算——只有方阵才值得先算行列式，非方阵一律直接消元。第 3 条②则是"给解求参数"类题的突破口：从"解不唯一"直接读出"系数矩阵奇异"。
\end{fdeep}
